% Part VIII — Head-to-Head and Why JuGeo Is Easier (Slides 141–160)

%% Slide 141 — Section divider
\begin{frame}
\frametitle{}
\vfill
\begin{center}
  {\Huge\color{jugeoBlue}\textbf{Part VI}}\\[12pt]
  {\LARGE\color{jugeoDark}The Same Proof --- Three Ways}\\[20pt]
  \begin{tikzpicture}
    \draw[jugeoBlue, line width=2pt] (-3,0) -- (3,0);
  \end{tikzpicture}\\[12pt]
  {\large\color{gray}\itshape Lean, F*, and JuGeo walk into a proof obligation\ldots}
\end{center}
\vfill
\end{frame}

%% Slide 142 — The Challenge
\begin{frame}[fragile]
\frametitle{The Challenge}
\begin{block}{Prove this function correct}
Given \texttt{values}, \texttt{bias}, \texttt{mod}, \texttt{keep}, compute:
\[
  \sum_{\substack{v \,\in\, \texttt{values} \\ v \bmod \texttt{mod}\,=\,\texttt{keep}}}
  \!\!\!(v + \texttt{bias})
\]
\end{block}
\begin{lstlisting}
def solve(values, bias, mod, keep):
    total = 0
    for v in values:
        if int(v) % mod == keep:
            total += int(v) + bias
    return total
\end{lstlisting}
\medskip
\centering
{\large Three languages. Three proof styles. One verdict.}
\end{frame}

%% Slide 143 — In LEAN 4 — Step 1
\begin{frame}[fragile]
\frametitle{In Lean~4 --- Step 1: Rewrite Your Code}
\begin{alertblock}{New language required}
You must rewrite the Python in Lean's functional style.
\end{alertblock}
\begin{lstlisting}[language=ML,morekeywords={def,where,List,Int,filter,foldl}]
def spec (values : List Int) (bias mod keep : Int) : Int :=
  (values.filter (fun v => v % mod == keep)).foldl
    (fun acc v => acc + v + bias) 0
\end{lstlisting}
\vfill
\begin{itemize}
  \item Learn Lean syntax, dependent types, universe levels
  \item Maintain \jugmark{two codebases}: Python + Lean
  \item No guarantee the Lean code matches the Python
\end{itemize}
\end{frame}

%% Slide 144 — In LEAN 4 — Step 2
\begin{frame}[fragile]
\frametitle{In Lean~4 --- Step 2: Write a Tactic Proof}
\begin{lstlisting}[language=ML,morekeywords={theorem,by,unfold,induction,with,nil,cons,simp,split,omega}]
theorem solve_eq_spec
    (values : List Int) (bias mod keep : Int)
    : solve values bias mod keep
      = spec values bias mod keep := by
  unfold solve spec
  induction values with
  | nil => simp
  | cons v vs ih => simp; split; omega
\end{lstlisting}
\vfill
\begin{columns}[T]
\begin{column}{0.48\textwidth}
\begin{alertblock}{Fragile}
Change one line of code $\Rightarrow$\\
proof breaks entirely
\end{alertblock}
\end{column}
\begin{column}{0.48\textwidth}
\begin{alertblock}{Manual}
Every branch, every case,\\
written by hand
\end{alertblock}
\end{column}
\end{columns}
\end{frame}

%% Slide 145 — In LEAN 4 — Step 3
\begin{frame}
\frametitle{In Lean~4 --- Step 3: Extract to C}
\begin{center}
\begin{tikzpicture}[
  box/.style={draw=jugeoBlue, fill=jugeoLight!30, rounded corners=4pt,
              minimum width=3cm, minimum height=1.2cm, align=center,
              font=\small\bfseries},
  arr/.style={-{Stealth[length=6pt]}, thick, jugeoBlue}
]
  \node[box] (lean) at (0,0) {Lean 4\\source};
  \node[box] (c) at (5,0) {C binary};
  \node[box, draw=jugeoRed, fill=jugeoRed!10] (py) at (5,-2.2)
    {Python?\\{\normalfont\color{jugeoRed}\itshape Not supported}};
  \draw[arr] (lean) -- node[above,font=\small] {extract} (c);
  \draw[arr, dashed, jugeoRed] (lean) -- (py);
\end{tikzpicture}
\end{center}
\vfill
\begin{itemize}
  \item Output carries \jugmark{no proof trace}
  \item If you want Python --- tough luck
  \item Proof checked at compile time, invisible at runtime
\end{itemize}
\end{frame}

%% Slide 146 — In F* — Step 1
\begin{frame}[fragile]
\frametitle{In F* --- Step 1: Rewrite Your Code}
\begin{alertblock}{Another new language}
F*'s ML-like syntax with totality annotations.
\end{alertblock}
\begin{lstlisting}[language=ML,morekeywords={let,rec,Tot,int,list,decreases,match,with,if,then,else}]
let rec solve (values: list int) (bias mod_ keep: int)
  : Tot int (decreases values) =
  match values with
  | [] -> 0
  | v :: vs ->
    if v % mod_ = keep
    then (v + bias) + solve vs bias mod_ keep
    else solve vs bias mod_ keep
\end{lstlisting}
\vfill
\begin{itemize}
  \item Must learn F* syntax \emph{and} totality annotations
  \item Still maintaining \jugmark{two codebases}
\end{itemize}
\end{frame}

%% Slide 147 — In F* — Step 2
\begin{frame}[fragile]
\frametitle{In F* --- Step 2: Write a Recursive Lemma}
\begin{lstlisting}[language=ML,morekeywords={let,rec,Lemma,ensures,decreases,match,with,list,int}]
let rec solve_correct (values: list int)
    (bias mod_ keep: int)
  : Lemma (ensures solve values bias mod_ keep
           = expected_total values bias mod_ keep)
          (decreases values) =
  match values with
  | [] -> ()
  | _ :: vs -> solve_correct vs bias mod_ keep
\end{lstlisting}
\vfill
\begin{columns}[T]
\begin{column}{0.48\textwidth}
\begin{exampleblock}{Pro}
Z3 auto-discharges arithmetic ---\\
less manual than Lean!
\end{exampleblock}
\end{column}
\begin{column}{0.48\textwidth}
\begin{alertblock}{Con}
When Z3 times out:\\[2pt]
\texttt{\color{jugeoRed}"Unknown."}\\[2pt]
No explanation. No fix.
\end{alertblock}
\end{column}
\end{columns}
\end{frame}

%% Slide 148 — In F* — Step 3
\begin{frame}
\frametitle{In F* --- Step 3: Extract to OCaml}
\begin{center}
\begin{tikzpicture}[
  box/.style={draw=jugeoBlue, fill=jugeoLight!30, rounded corners=4pt,
              minimum width=3cm, minimum height=1.2cm, align=center,
              font=\small\bfseries},
  arr/.style={-{Stealth[length=6pt]}, thick, jugeoBlue}
]
  \node[box] (fstar) at (0,0) {F*\\source};
  \node[box] (ocaml) at (5,0) {OCaml\\binary};
  \node[box, draw=jugeoRed, fill=jugeoRed!10] (py) at (5,-2.2)
    {Python?\\{\normalfont\color{jugeoRed}\itshape Not supported}};
  \draw[arr] (fstar) -- node[above,font=\small] {extract} (ocaml);
  \draw[arr, dashed, jugeoRed] (fstar) -- (py);
\end{tikzpicture}
\end{center}
\vfill
\begin{itemize}
  \item No provenance in extracted output
  \item No Python target
  \item Same pattern: proof stays behind, code ships alone
\end{itemize}
\end{frame}

%% Slide 149 — In JuGeo — Step 1
\begin{frame}[fragile]
\frametitle{In JuGeo --- Step 1: Write Python}
\begin{exampleblock}{No new language. This is your real code.}
\end{exampleblock}
\begin{lstlisting}
def solve(values, bias, mod, keep):
    total = 0
    for value in values:
        if int(value) % mod == keep:
            total += int(value) + bias
    return total
\end{lstlisting}
\vfill
\begin{center}
\begin{tikzpicture}
  \node[draw=jugeoGreen, fill=jugeoGreen!10, rounded corners=4pt,
        inner sep=10pt, font=\large\bfseries\color{jugeoGreen}]
    {That's it. Step 1 is done.};
\end{tikzpicture}
\end{center}
\end{frame}

%% Slide 150 — In JuGeo — Step 2
\begin{frame}[fragile]
\frametitle{In JuGeo --- Step 2: Write a Spec (Also Python)}
\begin{lstlisting}
def spec(result, values, bias, mod, keep):
    return isinstance(result, int) and \
           result == expected_total(values, bias, mod, keep)
\end{lstlisting}
\vfill
\begin{lstlisting}
def expected_total(values, bias, mod, keep):
    return sum(int(v) + bias
               for v in values
               if int(v) % mod == keep)
\end{lstlisting}
\vfill
\begin{itemize}
  \item Spec is a \jugmark{plain Python predicate}
  \item Returns \texttt{True} when \texttt{result} is correct
  \item No tactics, no lemmas, no new syntax
\end{itemize}
\end{frame}

%% Slide 151 — In JuGeo — Step 3
\begin{frame}[fragile]
\frametitle{In JuGeo --- Step 3: Declare a Cover}
\begin{lstlisting}
cover = [
    {"values": [],       "bias": 0, "mod": 2, "keep": 0},
    {"values": [4,7,10], "bias": 1, "mod": 3, "keep": 1},
    {"values": [1,2,3],  "bias": 5, "mod": 1, "keep": 0},
    # ... more points to span the input space
]
\end{lstlisting}
\vfill
\begin{center}
\begin{tikzpicture}[
  sbox/.style={draw=jugeoBlue, fill=jugeoLight!30, rounded corners=4pt,
               minimum width=3.4cm, minimum height=0.9cm, align=center,
               font=\small},
  arr/.style={-{Stealth[length=5pt]}, thick, jugeoBlue}
]
  \node[sbox] (cover) at (0,0) {Cover points};
  \node[sbox] (z3) at (4.5,0) {Z3 arithmetic};
  \node[sbox] (rt) at (9,0) {Runtime checks};
  \draw[arr] (cover) -- (z3);
  \draw[arr] (z3) -- (rt);
\end{tikzpicture}
\end{center}
\medskip
\centering JuGeo does the rest: Z3 for arithmetic, runtime for edge cases,\\
scaffold for \jugmark{provenance}.
\end{frame}

%% Slide 152 — In JuGeo — The Result
\begin{frame}
\frametitle{In JuGeo --- The Result}
\begin{center}
\begin{tikzpicture}[
  box/.style={draw=jugeoBlue, fill=jugeoLight!30, rounded corners=6pt,
              minimum width=3.2cm, minimum height=1.4cm, align=center,
              font=\small\bfseries},
  result/.style={draw=jugeoGreen, fill=jugeoGreen!10, rounded corners=6pt,
                 minimum width=5cm, minimum height=1.6cm, align=center,
                 font=\normalsize\bfseries, line width=1.5pt},
  arr/.style={-{Stealth[length=6pt]}, thick, jugeoBlue}
]
  \node[box] (code) at (-3.5,1.5) {Your Python\\code};
  \node[box] (spec) at (0,1.5) {Your Python\\spec};
  \node[box] (cover) at (3.5,1.5) {Your\\cover};
  \node[result, color=jugeoGreen] (out) at (0,-1)
    {Proof-carrying Python\\with multi-channel evidence};
  \draw[arr] (code) -- (out);
  \draw[arr] (spec) -- (out);
  \draw[arr] (cover) -- (out);
\end{tikzpicture}
\end{center}
\vfill
\begin{itemize}
  \item The \jugmark{same code} you wrote, now annotated
  \item Verification metadata travels with the artifact
  \item Any downstream verifier can re-check --- in Python
\end{itemize}
\end{frame}

%% Slide 153 — Side-by-Side Summary Table
\begin{frame}
\frametitle{Side-by-Side Summary}
\centering
\small
\renewcommand{\arraystretch}{1.35}
\begin{tabular}{@{}l
    >{\centering\arraybackslash}p{2.4cm}
    >{\centering\arraybackslash}p{2.4cm}
    >{\centering\arraybackslash}p{2.8cm}@{}}
\toprule
 & \textbf{Lean~4} & \textbf{F*} & \jugmark{JuGeo} \\
\midrule
\textbf{Write in}
  & Lean & F* & \color{jugeoGreen}Python \\
\textbf{Ship as}
  & C & OCaml & \color{jugeoGreen}Python (same!) \\
\textbf{Proof effort}
  & Manual tactics & Recursive lemma & \color{jugeoGreen}Cover + spec \\
\textbf{On failure}
  & \color{jugeoRed}``tactic failed'' & \color{jugeoRed}``Unknown'' & \color{jugeoGreen}Structured obstruction \\
\textbf{Trust model}
  & Binary & Binary & \color{jugeoGreen}7-level algebra \\
\bottomrule
\end{tabular}
\end{frame}

%% Slide 154 — Reason 1
\begin{frame}
\frametitle{Why JuGeo Is Easier --- Reason 1}
\vfill
\begin{center}
\begin{tikzpicture}
  \node[draw=jugeoGreen, fill=jugeoGreen!8, rounded corners=8pt,
        inner sep=16pt, text width=10cm, align=center,
        font=\Large\bfseries\color{jugeoDark}]
    {You don't rewrite your code.};
\end{tikzpicture}
\end{center}
\vfill
\begin{itemize}
  \item \jugmark{No second language} to learn
  \item No two codebases to keep in sync
  \item The code under verification \emph{is} the code you ship
\end{itemize}
\vfill
\end{frame}

%% Slide 155 — Reason 2
\begin{frame}
\frametitle{Why JuGeo Is Easier --- Reason 2}
\vfill
\begin{center}
\begin{tikzpicture}
  \node[draw=jugeoGreen, fill=jugeoGreen!8, rounded corners=8pt,
        inner sep=16pt, text width=10cm, align=center,
        font=\Large\bfseries\color{jugeoDark}]
    {You don't write tactic proofs.};
\end{tikzpicture}
\end{center}
\vfill
\begin{itemize}
  \item Declare the \jugmark{cover} and the \jugmark{spec}
  \item The semantic-move engine does the proof search
  \item Arithmetic goes to Z3; edge cases go to runtime
  \item You describe \emph{what} to verify, not \emph{how} to prove it
\end{itemize}
\vfill
\end{frame}

%% Slide 156 — Reason 3
\begin{frame}
\frametitle{Why JuGeo Is Easier --- Reason 3}
\vfill
\begin{center}
\begin{tikzpicture}
  \node[draw=jugeoGreen, fill=jugeoGreen!8, rounded corners=8pt,
        inner sep=16pt, text width=10cm, align=center,
        font=\Large\bfseries\color{jugeoDark}]
    {Failure is actionable.};
\end{tikzpicture}
\end{center}
\vfill
\begin{columns}[T]
\begin{column}{0.48\textwidth}
\begin{alertblock}{Lean / F*}
\texttt{tactic failed}\\
\texttt{Unknown}\\[4pt]
\emph{Now what?}
\end{alertblock}
\end{column}
\begin{column}{0.48\textwidth}
\begin{exampleblock}{JuGeo}
Structured \jugmark{obstruction}\\
with a \jugmark{repair frontier}:\\[4pt]
\emph{what} failed and \emph{how} to fix it
\end{exampleblock}
\end{column}
\end{columns}
\vfill
\end{frame}

%% Slide 157 — Reason 4
\begin{frame}
\frametitle{Why JuGeo Is Easier --- Reason 4}
\vfill
\begin{center}
\begin{tikzpicture}
  \node[draw=jugeoGreen, fill=jugeoGreen!8, rounded corners=8pt,
        inner sep=16pt, text width=10cm, align=center,
        font=\Large\bfseries\color{jugeoDark}]
    {Mixed evidence is natural.};
\end{tikzpicture}
\end{center}
\vfill
\begin{center}
\begin{tikzpicture}[
  chan/.style={draw=jugeoBlue, fill=jugeoLight!30, rounded corners=4pt,
              minimum width=3.2cm, minimum height=1cm, align=center, font=\small}
]
  \node[chan] (z3) at (0,0) {Z3 proved \textbf{8/10} points};
  \node[chan] (rt) at (5.5,0) {Tests covered \textbf{2} edge cases};
  \node[draw=jugeoGreen, fill=jugeoGreen!10, rounded corners=4pt,
        minimum width=5cm, minimum height=1cm, align=center, font=\small,
        line width=1.2pt] (sum) at (2.75,-1.5)
    {\jugmark{Honest account} --- not all-or-nothing};
  \draw[-{Stealth[length=5pt]}, thick, jugeoBlue] (z3) -- (sum);
  \draw[-{Stealth[length=5pt]}, thick, jugeoBlue] (rt) -- (sum);
\end{tikzpicture}
\end{center}
\vfill
\end{frame}

%% Slide 158 — Reason 5
\begin{frame}
\frametitle{Why JuGeo Is Easier --- Reason 5}
\vfill
\begin{center}
\begin{tikzpicture}
  \node[draw=jugeoGreen, fill=jugeoGreen!8, rounded corners=8pt,
        inner sep=16pt, text width=10cm, align=center,
        font=\Large\bfseries\color{jugeoDark}]
    {The proof ships with the code.};
\end{tikzpicture}
\end{center}
\vfill
\begin{center}
\begin{tikzpicture}[
  box/.style={draw=jugeoBlue, fill=jugeoLight!30, rounded corners=4pt,
              minimum width=2.6cm, minimum height=1cm, align=center, font=\small},
  arr/.style={-{Stealth[length=5pt]}, thick, jugeoGreen}
]
  \node[box] (dev) at (0,0) {Developer};
  \node[box, draw=jugeoGreen, fill=jugeoGreen!10] (art) at (4,0)
    {Verified\\artifact};
  \node[box] (ops) at (8,0) {Verifier};
  \draw[arr] (dev) -- node[above,font=\scriptsize]{deploy} (art);
  \draw[arr] (art) -- node[above,font=\scriptsize]{re-check} (ops);
\end{tikzpicture}
\end{center}
\medskip
\begin{itemize}
  \item Deployed Python artifact is \jugmark{self-describing}
  \item Any verifier can re-check without leaving Python
  \item No separate proof artifact to lose or mismatch
\end{itemize}
\vfill
\end{frame}

%% Slide 159 — What F* Can't Express
\begin{frame}[fragile]
\frametitle{What F* Can't Express --- Context Managers + Async}
\begin{lstlisting}
async def fetch(url):
    async with aiohttp.ClientSession() as session:
        async with session.get(url) as response:
            return await response.json()
\end{lstlisting}
\vfill
\begin{columns}[T]
\begin{column}{0.48\textwidth}
\begin{alertblock}{F* would need}
\begin{itemize}
  \item Custom effect for \texttt{aiohttp}
  \item Model for context managers
  \item Model for \texttt{async}/\texttt{await}
  \item Bindings for every library
\end{itemize}
\end{alertblock}
\end{column}
\begin{column}{0.48\textwidth}
\begin{exampleblock}{JuGeo}
\begin{itemize}
  \item Models it \jugmark{natively}
  \item Python is the language
  \item Runtime channel covers\\actual I/O behaviour
  \item No bindings needed
\end{itemize}
\end{exampleblock}
\end{column}
\end{columns}
\end{frame}

%% Slide 160 — The Bottom Line
\begin{frame}
\frametitle{The Bottom Line}
\vfill
\begin{center}
\begin{tikzpicture}
  \node[draw=jugeoBlue, fill=jugeoLight!30, rounded corners=10pt,
        inner sep=18pt, text width=11cm, align=center,
        line width=1.5pt] {
    {\LARGE\color{jugeoDark}\bfseries
      Write Python.\quad Get verified Python.\quad Ship verified Python.}\\[14pt]
    {\large\color{gray}\itshape
      All in one language, with honest trust accounting.}
  };
\end{tikzpicture}
\end{center}
\vfill
\begin{center}
\begin{tikzpicture}[
  lang/.style={draw=jugeoBlue, fill=jugeoLight!30, rounded corners=4pt,
               minimum width=2.2cm, minimum height=0.8cm, align=center,
               font=\small\bfseries},
  arr/.style={-{Stealth[length=5pt]}, thick, jugeoGreen}
]
  \node[lang] (w) at (0,0) {Write};
  \node[lang] (v) at (3.2,0) {Verify};
  \node[lang] (s) at (6.4,0) {Ship};
  \node[font=\scriptsize\color{jugeoGreen}, below=2pt] at (w.south) {Python};
  \node[font=\scriptsize\color{jugeoGreen}, below=2pt] at (v.south) {Python};
  \node[font=\scriptsize\color{jugeoGreen}, below=2pt] at (s.south) {Python};
  \draw[arr] (w) -- (v);
  \draw[arr] (v) -- (s);
\end{tikzpicture}
\end{center}
\vfill
\end{frame}
